\documentclass[a4paper]{article}

\usepackage{graphicx}
\usepackage{amsmath,amssymb}
\usepackage{minted}
\usepackage{multirow}
\usepackage{float}
\usepackage{algorithm}
\usepackage{algpseudocode}
\usepackage{todonotes}
\usepackage{forest}
\usepackage{xstring}
\usepackage{xcolor}
\usepackage[most]{tcolorbox}
\usepackage{xparse}
\usepackage{natbib}

\usetikzlibrary{graphs,graphdrawing}
\usegdlibrary{layered}

% for external graphviz
\input{/home/donkarlo/Dropbox/repo/nd_language_ai_project/src/nd_language_ai/formal/computer/programming/tex/latex/code_clips/external_graphviz.tex}

% for tags
\input{/home/donkarlo/Dropbox/repo/nd_language_ai_project/src/nd_language_ai/formal/computer/programming/tex/latex/part/control_sequence/kind/semtag/semtag.tex}

% for links
\input{/home/donkarlo/Dropbox/repo/nd_language_ai_project/src/nd_language_ai/formal/computer/programming/tex/latex/code_clips/seehere.tex}

\title{Ctrl+Shift+D Current-Line Deletion Fix}
\date{2026-09-04}
\author{Mohammad Rahmani}

\begin{document}
   \maketitle

   \section{Bug}
   Ctrl+Shift+D had been documented as deleting the current line, but the
   Source editor only implemented Ctrl+D for line duplication and the Visual
   editor had no Ctrl+Shift+D deletion handler.

   \section{Fix}
   Ctrl+Shift+D now deletes the logical line containing the active cursor in
   Source mode. The operation is one QTextCursor edit block, so Undo restores
   the removed line and no file-level delete operation is involved.

   In Visual mode, Ctrl+Shift+D deletes the active Visual block and preserves
   the metadata and block format of the following projected block. This keeps
   LaTeX serialization valid after deletion, including windowed large-document
   Visual editing.

   \bibliographystyle{plainnat}
   \bibliography{/home/donkarlo/Dropbox/repo/nd_shared_research/bibliography/bibliography.bib}
\end{document}
